\textbf{Área de la región}
\nopagebreak

La región está limitada por la curva \(y = \tan(x)\), el eje \(x\) y las rectas
\(x = 0\) y \(x = \frac{\pi}{3}\).  

Dividimos el intervalo \([0, \frac{\pi}{3}]\) en \(n\) partes iguales:
\[
\Delta x = \frac{\frac{\pi}{3} - 0}{n} = \frac{\pi}{3n}.
\]

Tomando puntos derechos:
\[
x_i = \frac{i \pi}{3n}.
\]

El área se expresa como límite de sumas de Riemann:
\[
A = \lim_{n \to \infty} \sum_{i=1}^{n} \tan\left(\frac{i \pi}{3n}\right) \Delta x.
\]

Sustituyendo \(\Delta x = \frac{\pi}{3n}\):
\[
A = \lim_{n \to \infty} \sum_{i=1}^{n} \tan\left(\frac{i \pi}{3n}\right) \cdot \frac{\pi}{3n}.
\]

Esta es la expresión de la integral como suma de Riemann. Evaluando la integral:

\[
A = \int_0^{\frac{\pi}{3}} \tan(x) \, dx
= \left[-\ln|\cos(x)|\right]_0^{\frac{\pi}{3}}.
\]

Sustituyendo los límites de integración:
\[
A = -\ln|\cos(\frac{\pi}{3})| + \ln|\cos(0)|
= -\ln\left(\frac{1}{2}\right) + \ln(1)
= \ln(2).
\]

Finalmente:
\[
\boxed{A = \ln(2)}.
\]
